% =============================================================================
% Project : THE_PHILOSOPHY_OF_VIABILITY_LIBRARY_PROJECT
% Book    : Book I — Ontology of Viability
% Book ID : VS.B01
% Part    : Part I — Identity and Manifestation
% Chapter : Chapter 08 — Comparative Philosophical Examination
% Chapter ID : VS.B01.P01.C08
%
% File    :
%   2_MAINMATTER/PART_I_IDENTITY_AND_MANIFESTATION/
%   CHAPTER_08_COMPARATIVE_PHILOSOPHICAL_EXAMINATION/
%   sections/section_05_aristotle.tex
%
% Section / Chapter / Asset:
%   Section 05 — Aristotle
%
% Structural Identifier :
%   VS.B01.P01.C08.S05
%
% Purpose :
%   To reconstruct Aristotle's treatment of being, substance, form,
%   matter, unity, potentiality, actuality, and change; examine the
%   relation between his account of substantial determination and
%   the problems of identity and manifestation; distinguish his
%   historical commitments from the foundations of the Philosophy
%   of Viability; and establish a critical transition to Plotinus.
%
% =============================================================================

\section{Aristotle}
\label{sec:ch08-aristotle}

\index[name]{Aristotle}
\index[concept]{Substance}
\index[concept]{Ousia}
\index[concept]{Form}
\index[concept]{Matter}
\index[concept]{Hylomorphism}
\index[concept]{Potentiality}
\index[concept]{Actuality}
\index[concept]{Being qua being}
\index[concept]{Pros hen}
\index[concept]{Substantial change}
\index[concept]{Principle of non-contradiction}

% =============================================================================
% 1. ARISTOTLE'S ONTOLOGICAL PROJECT AND ITS TEXTUAL STRUCTURE
% =============================================================================

\subsection{Aristotle's Ontological Project and Its Textual Structure}

\phantomsection\label{para:VS.B01.P01.C08.S05.SS01.P001}

Aristotle's significance for the present comparative examination
lies in the systematic connection he establishes among being,
substance, form, matter, actuality, potentiality, and change.
His investigations do not constitute a single treatise organized
according to the architecture of a modern philosophical system.
The works conventionally grouped under the titles
\textit{Categories}, \textit{Physics}, and \textit{Metaphysics}
address related problems from distinguishable analytical
perspectives. Their relation must therefore be reconstructed
rather than assumed to express an entirely uniform terminology
or an unchanging theoretical position. \cite{AristotleCategories,
AristotlePhysics,AristotleMetaphysics}

\phantomsection\label{para:VS.B01.P01.C08.S05.SS01.P002}

The \textit{Categories} investigates the ways in which beings
are classified and predicated. The \textit{Physics} examines
natural beings insofar as they undergo motion and change.
The \textit{Metaphysics} investigates being in a more general
respect, while repeatedly returning to the question of substance
and its principles. These works overlap, but their questions
are not interchangeable. In particular, the account of primary
substance in the \textit{Categories} cannot simply be substituted
for the later investigation of matter, form, and the composite
in the central books of the \textit{Metaphysics}. \footnote{For
the relation among the principal questions and textual strata
of the \textit{Metaphysics}, see Cohen and Reeve's reconstruction
of Aristotle's inquiry into being and substance
\cite{CohenReeveAristotleMetaphysics2025}.}

\phantomsection\label{para:VS.B01.P01.C08.S05.SS01.P003}

The present examination concentrates on the ontological
consequences of these investigations. It asks how Aristotle
understands the unity and determination of a substance, how
a substance may undergo change, and what role form and matter
play in explaining the constitution and transformation of
concrete beings. The question is not whether Aristotle already
possesses the conceptual distinctions developed in the
Philosophy of Viability. It is whether his analyses disclose
problems, explanatory resources, or limitations that must
be confronted by an ontology of identity and manifestation.

\phantomsection\label{para:VS.B01.P01.C08.S05.SS01.P004}

The distinction between Aristotle's own claims and the
interpretive reconstruction offered here is essential.
Aristotle's terminology belongs to a historically situated
philosophical framework. Its translation into contemporary
terms may clarify a problem, but translation does not establish
conceptual equivalence. The following discussion therefore
retains the principal Greek terms where their meanings carry
philosophical weight and treats disagreements among interpreters
as genuine interpretive questions rather than as obstacles
to be removed by an artificially unified account.

% =============================================================================
% 2. BEING SAID IN MANY WAYS
% =============================================================================

\subsection{Being Said in Many Ways and the Unity of Ontological Inquiry}

\phantomsection\label{para:VS.B01.P01.C08.S05.SS02.P001}

One of Aristotle's most consequential formulations is that
being is said in many ways. The claim appears in the
\textit{Metaphysics}, where Aristotle rejects the expectation
that every use of the verb ``to be'' must possess one
univocal meaning. Substance, quality, quantity, relation,
and other categories are all spoken of as beings, but they
are not beings in precisely the same manner. The plurality
of senses does not, however, make the inquiry into being
impossible. Aristotle argues that the different senses are
related to a central case: substance. \footnote{See
\textit{Metaphysics} $\Gamma$.2, 1003a21--1003b19.
The expression commonly rendered as ``being is said in
many ways'' is connected with the relation described as
\textit{pros hen}, or reference to one focal sense.}

\phantomsection\label{para:VS.B01.P01.C08.S05.SS02.P002}

The relation is commonly described as focal or pros-hen
equivocity. Qualities, quantities, and relations are beings
because they are predicated of, or otherwise belong to,
substances. Aristotle's examples of the healthy and the
medical illustrate the structure: different things may be
called healthy because they produce, preserve, indicate,
or possess health. Their uses are not unrelated homonyms,
but neither do they share one identical definition. They
are ordered through their reference to a central case.
The corresponding account of being gives substance a
privileged explanatory position. \cite{AristotleMetaphysics,
CohenReeveAristotleMetaphysics2025}

\phantomsection\label{para:VS.B01.P01.C08.S05.SS02.P003}

This account offers a way to investigate ontological
plurality without reducing every kind of being to one
unqualified sense. It also raises a question about the
unity of the inquiry itself: whether the different senses
of being form a sufficiently ordered field for a single
science to investigate them. Aristotle's answer depends
upon the explanatory centrality of substance. The unity
of the inquiry is therefore not achieved by erasing
differences among categories, but by relating those
differences to a focal ontological question.

\phantomsection\label{para:VS.B01.P01.C08.S05.SS02.P004}

The relevance to the Philosophy of Viability is methodological
and comparative. An ontology may investigate several
irreducible modes or respects of being while still asking
whether their relations form an intelligible whole.
Yet Aristotle's focal relation must not be identified
with the proposed architecture of the realms and hubs.
The seven-domain proposal is not a classification of
Aristotelian categories, and no direct correspondence
between substance and any one of its domains follows
from the comparison. The question of how distinct
ontological determinations may belong to a coherent
reality remains independently open.

% =============================================================================
% 3. THE CATEGORIES AND PRIMARY SUBSTANCE
% =============================================================================

\subsection{The Categories and the Priority of Primary Substance}

\phantomsection\label{para:VS.B01.P01.C08.S05.SS03.P001}

In the \textit{Categories}, Aristotle distinguishes ten
categories of predication and being, including substance,
quantity, quality, relation, place, time, position, state,
action, and affection. The classification is not merely
a list of grammatical expressions. It concerns different
ways in which something may be said to be and different
kinds of predicates that may be attributed to a subject.
Substance occupies the first position because individual
substances function as fundamental subjects of predication
within this account. \footnote{See \textit{Categories}
1b25--2a10 for the classification of things said of a
subject and things present in a subject, and 4a10--21
for the account of the categories.}

\phantomsection\label{para:VS.B01.P01.C08.S05.SS03.P002}

Aristotle distinguishes primary substances from secondary
substances. A particular human being or a particular
horse is a primary substance. The species and genus
under which such an individual is classified are
secondary substances. Primary substances are not said
of a subject in the manner in which a species or genus
is said of an individual. They also do not exist in
a subject in the manner in which an accidental quality
or quantity exists in something else. The distinction
establishes a hierarchy of predication and dependence
within the account of the \textit{Categories}, not
a hierarchy of ontological domains in the sense used
by the Philosophy of Viability.

\phantomsection\label{para:VS.B01.P01.C08.S05.SS03.P003}

The priority of primary substance is expressed through
the dependence of other categorical items upon individual
substances. A quality such as whiteness is present in
a subject; a species such as human is said of particular
human beings. If primary substances did not exist,
the relevant dependent items would have no subjects
in which to be present and no individuals of which
they could be predicated. Aristotle's argument thus
connects the basic status of individual substances
with the conditions under which other beings in the
categories can be said to exist. \cite{AristotleCategories,
CohenReeveAristotleMetaphysics2025}

\phantomsection\label{para:VS.B01.P01.C08.S05.SS03.P004}

This account is philosophically important because it
distinguishes the individual subject from the predicates,
classifications, and qualities associated with it.
The distinction prevents a simple reduction of an
individual to the sum of the descriptions that can
be predicated of it. Nevertheless, the \textit{Categories}
does not yet provide the complete analysis of what
makes a substance the substance it is. Its account
of primary substance identifies fundamental subjects
within a scheme of predication, while the later
investigation asks more directly what substance is
and which principles explain its constitution.

\phantomsection\label{para:VS.B01.P01.C08.S05.SS03.P005}

The distinction between primary and secondary substance
also requires interpretive care. The status of species
and genera, the relation between predication and
ontological dependence, and the relation between
the \textit{Categories} and the central books of
the \textit{Metaphysics} have generated sustained
scholarly debate. The present comparison therefore
treats the \textit{Categories} as a significant
account of individual substances and categorical
dependence, not as an exhaustive statement of every
Aristotelian commitment concerning substance. \footnote{For
the interpretive questions surrounding primary substances
and the relation between the \textit{Categories} and
the later hylomorphic analysis, see
\cite{CohenReeveAristotleMetaphysics2025}.}

% =============================================================================
% 4. SUBSTANCE, ESSENCE, AND THE SUBJECT
% =============================================================================

\subsection{Substance, Essence, and the Question of What a Thing Is}

\phantomsection\label{para:VS.B01.P01.C08.S05.SS04.P001}

In the central books of the \textit{Metaphysics}, Aristotle
reopens the question of substance. The inquiry is no
longer limited to identifying the subjects of predication.
It asks what substance is and what makes a substance
the kind of being it is. Aristotle presents several
candidates, including the subject, essence, universal,
and genus. The investigation does not simply repeat
the account of primary substance in the \textit{Categories};
it examines the principles that explain substantial
being. \footnote{See \textit{Metaphysics} Z.1--3,
1028b2--1029a30.}

\phantomsection\label{para:VS.B01.P01.C08.S05.SS04.P002}

The subject criterion initially appears attractive.
A subject is that of which other things are predicated
while it is not itself predicated of another subject.
Yet, when this criterion is applied to the analysis
of a concrete substance, it generates a difficulty.
Matter, form, and the composite of matter and form
may each be considered a subject in a relevant sense.
If every predicate is removed from a composite, the
remaining subject appears to be matter stripped of
all determination. Such matter, however, does not
sufficiently explain why the substance is a determinate
individual of a particular kind.

\phantomsection\label{para:VS.B01.P01.C08.S05.SS04.P003}

Aristotle consequently gives substantial importance
to essence and form. Essence concerns what a thing
is in itself, rather than an accidental feature that
may belong to it or cease to belong to it. The distinction
between essential and accidental predication allows
Aristotle to investigate why some predicates belong
to a thing in virtue of what it is, while others belong
to it contingently. The question is not exhausted
by identifying the subject that receives predicates;
it also concerns the explanatory account of the
subject's being the kind of thing it is. \cite{AristotleMetaphysics}

\phantomsection\label{para:VS.B01.P01.C08.S05.SS04.P004}

The analysis of essence is connected with definition,
but the relation between a definition and the being
of the defined substance must be handled carefully.
Aristotle does not merely offer a theory of the meanings
of words. His concern is with what a thing is and
with the account that expresses its essence. Yet the
status of universals, species, and substantial forms
remains contested in the interpretation of the
\textit{Metaphysics}. It is therefore unsafe to assume
that every essence is a universal entity existing
independently of the particular substances that
instantiate or express it. \footnote{For the discussion
of essence and definition, see \textit{Metaphysics}
Z.4--6, especially 1029b--1032a. The interpretation
of substantial form and its relation to universality
remains disputed; see \cite{CohenReeveAristotleMetaphysics2025}.}

\phantomsection\label{para:VS.B01.P01.C08.S05.SS04.P005}

Aristotle's analysis consequently advances beyond
a purely predicative account of substance. It seeks
an explanatory principle through which the determinate
being of a substance can be understood. This is a
substantial contribution to the ontology of determination:
the account of a thing must address not only what
can be said of it, but also the principle in virtue
of which it is the determinate being that it is.
The contribution remains specific to Aristotle's
framework and does not establish that essence,
form, or definition is identical with Identity.

% =============================================================================
% 5. MATTER, FORM, AND THE COMPOSITE
% =============================================================================

\subsection{Matter, Form, and the Constitution of Concrete Substance}

\phantomsection\label{para:VS.B01.P01.C08.S05.SS05.P001}

Aristotle's analysis of natural substances introduces
a distinction between matter (\textit{hylē}) and form
(\textit{eidos} or \textit{morphē}). Matter is that
from which a thing is constituted or that which serves
as an underlying subject in a process of change.
Form is the principle through which the matter is
determined as a substance of a particular kind.
The distinction is not simply between physical
material and an externally imposed shape. In natural
substances, form concerns the organization and
principle of being through which a living or natural
thing is the kind of thing it is. \cite{AristotlePhysics,
AristotleMetaphysics}

\phantomsection\label{para:VS.B01.P01.C08.S05.SS05.P002}

The relation is often called hylomorphism, from the
Greek terms for matter and form. A bronze statue,
for example, may be analysed in terms of bronze
as matter and the statue's form as the organization
through which that matter constitutes a statue.
The example is useful, but natural substances cannot
be reduced to artefacts. In a living organism, form
is not merely an external outline imposed by a maker.
It is associated with the organized capacities and
activities through which the organism lives and
operates. The analogy with artefacts must therefore
remain limited.

\phantomsection\label{para:VS.B01.P01.C08.S05.SS05.P003}

The composite of matter and form is central to the
analysis of sensible substances. Matter considered
in abstraction from form does not adequately explain
the determinate being of a concrete substance.
Form, in turn, is not simply an additional property
placed beside other properties. It has an explanatory
role in accounting for the substance's constitution,
organization, and characteristic activities.
Aristotle's account therefore resists both a reduction
of the substance to undifferentiated matter and a
description of its form as an arbitrary external
addition. \footnote{See \textit{Physics} I.7,
190b10--191a22, and \textit{Metaphysics} Z.3,
1029a1--30; H.1--6, 1042a24--1045b23.}

\phantomsection\label{para:VS.B01.P01.C08.S05.SS05.P004}

The distinction also clarifies why the same material
constituents may enter into different organized
substances and why the constitution of a substance
cannot be explained by listing its material components
alone. The explanatory question concerns the relation
between what a thing is made from and the organization
through which those materials constitute that thing.
The relation is not a simple identity between matter
and substance, nor does the presence of a particular
form automatically settle every question about the
numerical identity of the composite.

\phantomsection\label{para:VS.B01.P01.C08.S05.SS05.P005}

The following table summarizes the principal roles
of matter, form, and composite within the Aristotelian
analysis. It is an interpretive aid, not a claim that
these terms have one identical function in every
Aristotelian context.

% =============================================================================
% Project : THE_PHILOSOPHY_OF_VIABILITY_LIBRARY_PROJECT
% Book    : Book I — Ontology of Viability
% Book ID : VS.B01
% Part    : Part I — Identity and Manifestation
% Chapter : Chapter 08 — Comparative Philosophical Examination
% Chapter ID : VS.B01.P01.C08
%
% File :
%   4_ASSETS/tables/VS.B01.P01.C08.S05.SS05.P005.TL01.aristotelian_matter_form_composite.tex
%
% Section / Chapter / Asset:
%   Aristotelian Matter, Form, and Composite
%
% Structural Identifier :
%   VS.B01.P01.C08.S05.SS05.P005.TL01
%
% Asset Type : TL
%
% Purpose :
%   To distinguish the explanatory roles of matter, form, and
%   composite in Aristotle's hylomorphic analysis.
%
% =============================================================================

\begin{table}[htbp]
\centering
\caption{Matter, Form, and Composite in Aristotle's Hylomorphic Analysis}
\label{tab:aristotle-matter-form-composite}

\small
\renewcommand{\arraystretch}{1.3}

\begin{tabularx}{\textwidth}{
    >{\raggedright\arraybackslash}p{0.17\textwidth}
    >{\raggedright\arraybackslash}p{0.37\textwidth}
    >{\raggedright\arraybackslash}X
}
\toprule
\textbf{Aspect} &
\textbf{Explanatory Role} &
\textbf{Interpretive Limit}
\\
\midrule

Matter (\textit{hylē})
&
That from which a substance is constituted, or the
underlying subject considered in relation to a process
of change.
&
Matter considered without substantial determination
does not by itself explain why a concrete substance
is the kind of thing it is.
\\

\addlinespace

Form (\textit{eidos}/\textit{morphē})
&
The principle through which matter is determined as
a substance of a particular kind and through which
its characteristic organization is explained.
&
Form must not be reduced to an externally imposed
shape or treated as an independently established
equivalent of identity.
\\

\addlinespace

Composite
&
The concrete substance considered as constituted
through the relation of matter and form.
&
The composite is not merely an aggregate of
independently complete constituents; its unity
requires an account of substantial organization.
\\

\addlinespace

Potentiality and actuality
&
The relation through which matter and form may be
analysed in terms of capacity and realization.
&
This relation has several applications in Aristotle
and must not be reduced to one temporal sequence
or one kind of change.
\\

\bottomrule
\end{tabularx}

\end{table}

% =============================================================================
% END OF TABLE TL01
% =============================================================================

% =============================================================================
% 6. SUBSTANTIAL UNITY AND COMPOSITION
% =============================================================================

\subsection{Substantial Unity and the Problem of Composition}

\phantomsection\label{para:VS.B01.P01.C08.S05.SS06.P001}

The relation between matter and form raises a further
question: how can a composite be one substance rather
than an aggregate of separately constituted elements?
If matter and form were merely two independent things
placed alongside one another, their combination would
require an additional explanation. Aristotle's account
therefore treats form as a principle through which
the matter constitutes one determinate substance.
The unity of the composite is connected with its
substantial organization, not merely with the fact
that its components occupy one location.

\phantomsection\label{para:VS.B01.P01.C08.S05.SS06.P002}

In \textit{Metaphysics} H.6, Aristotle examines the
unity of matter and form and asks why a definition
or account of a substance is one rather than a mere
sequence of statements. His response is connected
with the relation between potentiality and actuality:
matter is potentially what the composite is, while
form is associated with the actuality through which
the substance is constituted. The explanation is
not a general theory of every kind of unity, but
an account of substantial unity within the hylomorphic
framework. \footnote{See \textit{Metaphysics} H.6,
1045a7--b23.}

\phantomsection\label{para:VS.B01.P01.C08.S05.SS06.P003}

This analysis provides an important alternative
to an account that treats every whole as an aggregate
of independently complete parts. A substantial
whole may possess a mode of unity that depends
upon the organization of its constituents and upon
the role those constituents perform within the
whole. Nevertheless, the distinction between a
substance and an aggregate cannot be settled by
the mere presence of organization. The relevant
account must specify what kind of organization is
involved and why it has substantial significance.

\phantomsection\label{para:VS.B01.P01.C08.S05.SS06.P004}

The comparison with the Philosophy of Viability
must preserve the distinction between Unity,
Composition, Aggregate, and Identity. Aristotle's
account of substantial unity may illuminate how
a plurality of material constituents can belong
to one organized substance. It does not establish
that every unity corresponds to one identity, that
every composite possesses exactly one identity,
or that the unity of a substance is identical with
its identity. The number and status of identities
remain questions requiring independent argument.

\phantomsection\label{para:VS.B01.P01.C08.S05.SS06.P005}

The Aristotelian account is therefore relevant
as a theory of substantial organization, not as
a ready-made solution to the problem of identity.
It provides an explanatory model in which the
unity of a natural substance is related to its
form and matter, while leaving open whether
the same analysis can be extended to every
kind of existent or to entities whose organization
does not conform to the structure of natural
substances.

% =============================================================================
% 7. POTENTIALITY AND ACTUALITY
% =============================================================================

\subsection{Potentiality and Actuality}

\phantomsection\label{para:VS.B01.P01.C08.S05.SS07.P001}

The distinction between potentiality (\textit{dynamis})
and actuality (\textit{energeia} or \textit{entelecheia})
extends Aristotle's analysis beyond the static
description of substance. Potentiality concerns
a capacity or principle through which something
may act, undergo change, or attain a determinate
condition. Actuality concerns the exercise or
realization of such a capacity. The distinction
is not limited to the possibility of physical
motion. Aristotle extends it to the relation
between matter and form, capacities and their
exercise, and different modes of being. \cite{AristotleMetaphysics}

\phantomsection\label{para:VS.B01.P01.C08.S05.SS07.P002}

Potentiality must not be understood as a merely
subjective expectation or as a possibility that
exists only because an observer can imagine it.
Aristotle's discussion concerns capacities and
principles belonging to beings. A material subject
may be potentially capable of receiving a form;
a living being may possess capacities that are
not currently exercised; and an agent may possess
the capacity to produce a change. These cases
are related, but they do not all instantiate
one simple meaning of potentiality.

\phantomsection\label{para:VS.B01.P01.C08.S05.SS07.P003}

Actuality is correspondingly broader than the
completion of a process of movement. Aristotle
uses examples such as seeing and having sight,
being awake and being asleep, and the relation
between a formed object and the matter from
which it is made. These examples show that
actuality and potentiality are analogically
related across different domains. Their meanings
cannot be reduced to a single temporal sequence
in which every potentiality must first exist
and then become an actuality of the same kind.
\footnote{See \textit{Metaphysics} $\Theta$.6,
1048a25--b9, and $\Theta$.8, 1049b4--1051a3.}

\phantomsection\label{para:VS.B01.P01.C08.S05.SS07.P004}

The distinction has explanatory significance
for the hylomorphic account. Matter is related
to form as potentiality to actuality, although
Aristotle's use of these concepts is not exhausted
by this relation. The form of a substance is
not merely a future state that matter may later
receive. In the completed substance, form is
associated with the actuality through which
the matter constitutes the determinate being
that it is. The relation therefore concerns
the structure of being, not only a sequence
of events.

\phantomsection\label{para:VS.B01.P01.C08.S05.SS07.P005}

The conceptual relation between potentiality
and actuality is represented schematically below.
The diagram distinguishes their explanatory
roles without presenting them as two independent
substances or as a universal linear hierarchy
of all beings.

% =============================================================================
% Project : THE_PHILOSOPHY_OF_VIABILITY
% Book    : Book I — Ontology of Viability
% Book ID : VS.B01
% Part    : Part I — Identity and Manifestation (VS.B01.P01)
% Chapter : Chapter 08 — Comparative Philosophical Examination
% Chapter ID : VS.B01.P01.C08
%
% File    : 4_ASSETS/diagrams/VS.B01.P01.C08.S05.SS07.P005.DG01.potentiality_actuality_and_form.tex
%
% Section / Chapter / Asset:
%   Diagram DG01 — Potentiality, Actuality, and Form
%
% Structural Identifier : VS.B01.P01.C08.S05.SS07.P005.DG01
% Asset Type : DG
%
% Purpose :
%   To display the explanatory relation among matter, potentiality, form, and
%   actuality without representing them as a universal linear hierarchy of
%   being.
%
% Representational limit :
%   The arrangement is explanatory rather than a universal sequence or
%   hierarchy of all beings.
%
% Dependencies :
%   TikZ; libraries used: arrows.meta, positioning. All are loaded by the
%   project preamble; no per-file package or library loading.
%
% =============================================================================

\begin{figure}[htbp]
    \centering
    \hyphenpenalty=10000 \exhyphenpenalty=10000\relax  % no hyphenation inside diagram labels

    \begin{tikzpicture}[
        font=\small,
        % --- base box ---
        box/.style={
            draw, rounded corners=3pt, align=center,
            text width=2.9cm, minimum height=1.5cm,
            inner sep=5pt, line width=0.7pt
        },
        % --- node roles ---
        concept/.style={box, fill=blue!8, draw=blue!55!black},
        glossbox/.style={box, dashed, fill=gray!10, draw=black!60},
        % --- arrows and labels ---
        directed/.style={-{Stealth[length=2.5mm]}, line width=0.9pt},
        vslink/.style={draw=black!55, dashed, line width=0.7pt},
        elabel/.style={font=\scriptsize\itshape, text=black!75, align=center}
    ]

    % ---------------------------------------------------------------
    % Explanatory arrangement: matter - potentiality - form - actuality
    % ---------------------------------------------------------------
    \node[concept] (matter) {
        \textbf{Matter}\\[2pt] \textit{hylē}
    };
    \node[concept, right=2.6cm of matter] (potentiality) {
        \textbf{Potentiality}\\[2pt] \textit{dynamis}
    };
    \node[concept, below=2.0cm of potentiality] (form) {
        \textbf{Form}\\[2pt] \textit{eidos / morphē}
    };
    \node[concept, right=2.6cm of form] (actuality) {
        \textbf{Actuality}\\[2pt] \textit{energeia / entelecheia}
    };

    % ---------------------------------------------------------------
    % Glosses (dashed, undirected)
    % ---------------------------------------------------------------
    \node[glossbox, below=1.2cm of matter] (matternote) {
        Underlying subject and capacity for determination
    };
    \node[glossbox, below=1.2cm of actuality] (actualitynote) {
        Realization and determinate constitution
    };
    \draw[vslink] (matter)    -- (matternote);
    \draw[vslink] (actuality) -- (actualitynote);

    % ---------------------------------------------------------------
    % Relations
    % ---------------------------------------------------------------
    \draw[directed] (matter) -- node[elabel, above=2pt, text width=2.4cm]
        {considered as capable} (potentiality);
    \draw[directed] (potentiality) -- node[elabel, right=4pt, text width=2.6cm]
        {relation to determination} (form);
    \draw[directed] (form) -- node[elabel, above=2pt, text width=2.4cm]
        {actualizing role} (actuality);

    \end{tikzpicture}

    \caption{A schematic presentation of the relation between matter,
    potentiality, form, and actuality in Aristotle's hylomorphic
    analysis. The arrangement is explanatory rather than a
    universal sequence or hierarchy of all beings.}
    \label{fig:aristotle-potentiality-actuality-form}
\end{figure}

% =============================================================================
% END OF DIAGRAM DG01
% =============================================================================


% =============================================================================
% 8. CHANGE, CONTINUITY, AND GENERATION
% =============================================================================

\subsection{Change, Continuity, and the Difference Between Alteration and Generation}

\phantomsection\label{para:VS.B01.P01.C08.S05.SS08.P001}

Aristotle's account of change is closely connected
with the distinction between matter, form, and
privation. In natural change, a subject undergoes
a transition from one condition to another.
The analysis distinguishes what persists through
the change from the determination that is acquired
or lost. In cases of alteration, a substance may
remain the same while receiving contrary qualities.
In cases of substantial generation and destruction,
the question is more difficult because the change
concerns the coming to be or passing away of a
substance itself. \footnote{See \textit{Physics}
I.7, 190a13--191a22, and II.1, 192b8--23.}

\phantomsection\label{para:VS.B01.P01.C08.S05.SS08.P002}

The distinction between accidental change and
substantial change prevents the assumption that
every alteration has the same ontological structure.
A particular human being may become pale or dark
without thereby becoming a different substance.
Such a change affects an accidental determination.
By contrast, generation and destruction concern
the coming to be and passing away of substances.
The underlying matter plays a different explanatory
role in these cases, since the form that constitutes
a substance is not simply an accidental predicate
that can be exchanged while the same substance
remains unchanged in kind.

\phantomsection\label{para:VS.B01.P01.C08.S05.SS08.P003}

Aristotle's account thus provides a differentiated
analysis of continuity. The persistence of a subject
through alteration is not identical with the
continuity involved in substantial generation.
Nor can the persistence of matter by itself
establish the persistence of the substance whose
matter it is. The analysis requires attention
to the kind of change, the subject that undergoes
it, and the determination that is acquired or lost.
These distinctions make it possible to examine
continuity without assuming that all changes
preserve the same ontological subject.

\phantomsection\label{para:VS.B01.P01.C08.S05.SS08.P004}

The Aristotelian distinction is especially relevant
to the problem of replacement. If matter persists
while one substantial form is lost and another
is acquired, the continued presence of the matter
does not by itself show that the same substance
has survived. Conversely, a change in accidental
qualities does not automatically entail the
destruction of the substance. The relation between
the persistence of a subject and the identity
of the substance therefore requires a more
precise account than the observation that
something remains through change.

\phantomsection\label{para:VS.B01.P01.C08.S05.SS08.P005}

The illustration below offers a concrete visual
analogy for the distinction between material
substrate, organization, and the resulting
composite. It is intended to support the reader's
imagination and to provide a visual pause in
the discussion; it is not evidence for the
Aristotelian theory and does not represent every
kind of natural substance.

% =============================================================================
% Project : THE_PHILOSOPHY_OF_VIABILITY_LIBRARY_PROJECT
% Book    : Book I — Ontology of Viability
% Book ID : VS.B01
% Part    : Part I — Identity and Manifestation
% Chapter : Chapter 08 — Comparative Philosophical Examination
% Chapter ID : VS.B01.P01.C08
%
% File :
%   4_ASSETS/illustrations/VS.B01.P01.C08.S05.SS08.P005.IL01.matter_form_and_artifact.tex
%
% Section / Chapter / Asset:
%   Matter, Form, and Artifact
%
% Structural Identifier :
%   VS.B01.P01.C08.S05.SS08.P005.IL01
%
% Asset Type : IL
%
% Purpose :
%   To provide an aesthetically composed visual analogy for the
%   distinction between material substrate, organization, and
%   the resulting artifact.
%
% Dependencies :
%   TikZ; positioning library, loaded by the project preamble.
%
% Representational limit :
%   The illustration concerns an artifact and must not be taken
%   as a complete representation of natural substances or as
%   evidence for hylomorphism.
%
% =============================================================================

\begin{figure}[htbp]
\centering

\begin{tikzpicture}[
    material/.style={
        draw=ViabilityGold,
        fill=ViabilityPaleGold,
        rounded corners=2mm,
        line width=0.8pt,
        minimum width=34mm,
        minimum height=17mm,
        align=center,
        text=ViabilityBlack
    },
    form/.style={
        draw=ViabilityDarkBlue,
        fill=ViabilityPaleBlue,
        rounded corners=2mm,
        line width=0.8pt,
        minimum width=34mm,
        minimum height=17mm,
        align=center,
        text=ViabilityBlack
    },
    result/.style={
        draw=ViabilityGreen,
        fill=ViabilityPaleGreen,
        rounded corners=2mm,
        line width=0.8pt,
        minimum width=43mm,
        minimum height=21mm,
        align=center,
        text=ViabilityBlack
    },
    flow/.style={
        -{Stealth[length=2.5mm]},
        line width=1pt,
        draw=ViabilityGray
    }
]

\node[material] (wood) {
    \textbf{Material}\\
    Wood
};

\node[form, right=17mm of wood] (organization) {
    \textbf{Organization}\\
    Shape and arrangement
};

\node[result, below=18mm of $(wood)!0.5!(organization)$] (artifact) {
    \textbf{Artifact}\\
    A constituted wooden object
};

\draw[flow] (wood.south) -- ++(0,-6mm) -| (artifact.west);
\draw[flow] (organization.south) -- ++(0,-6mm) -| (artifact.east);

\node[
    below=7mm of artifact,
    text=ViabilityGray,
    align=center,
    font=\small,
    text width=105mm
] {
    A limited analogy for the relation between what a thing is made
    from and the organization through which it constitutes an artifact.
};

\end{tikzpicture}

\caption{A visual analogy for the distinction between material
substrate and organization in an artifact. The illustration is
not a complete model of Aristotle's account of natural substances.}
\label{fig:aristotle-matter-form-artifact}

\end{figure}

% =============================================================================
% END OF ILLUSTRATION IL01
% =============================================================================

% =============================================================================
% 9. IDENTITY, DETERMINATION, AND THE LIMITS OF THE COMPARISON
% =============================================================================

\subsection{Identity, Determination, and the Limits of Aristotelian Explanation}

\phantomsection\label{para:VS.B01.P01.C08.S05.SS09.P001}

Aristotle's account of substance, essence, and form
raises a question central to the present book:
whether the principle that explains what a thing
is can also explain the identity of that thing.
The question is not answered merely by observing
that form is essential to the constitution of
a substance. An account of what makes something
a substance of a particular kind and an account
of the identity of that substance may be closely
related, but their equivalence requires an argument
that Aristotle's terminology cannot supply for
the Philosophy of Viability without further analysis.

\phantomsection\label{para:VS.B01.P01.C08.S05.SS09.P002}

The Aristotelian notion of essence concerns what
a substance is in itself. The present ontology,
however, does not define Identity as a collection
of essential predicates, as a species, or as
a substantial form. Identity is not an additional
property attached to an otherwise identity-neutral
being, but neither is it simply interchangeable
with Determination. The analytical distinction
between Being and Identity, their ontological
inseparability, and the absence of ontological
priority between them remain three independent
principles. Aristotle's analysis may illuminate
the problem of determination, but it does not
replace the argument required for these principles.

\phantomsection\label{para:VS.B01.P01.C08.S05.SS09.P003}

The distinction also matters for plurality.
Aristotle's account of a substantial individual
may provide reasons for treating a particular
natural substance as one organized being.
It does not follow that every organized whole
has exactly one identity, that every plurality
of parts constitutes a plurality of identities,
or that the unity of a substantial form settles
the numerical question of identity. The possibility
of multiple identities, including identities
considered across different ontological respects
or levels, remains an open problem in the
Philosophy of Viability.

\phantomsection\label{para:VS.B01.P01.C08.S05.SS09.P004}

Aristotle's account of actuality and potentiality
also cannot be equated with the account of
Invariance developed in this book. A capacity
may be actualized, a form may be acquired,
and a substance may undergo alteration without
these processes being described simply as
the preservation of an invariant. Any claim
of invariance requires a specified transformation,
domain, respect, type, and limit. Aristotle's
conceptual resources may help formulate questions
about what remains through change, but they
do not establish that actuality, form, or
substance is an invariant.

\phantomsection\label{para:VS.B01.P01.C08.S05.SS09.P005}

The principle of non-contradiction provides
another point of comparison. Aristotle formulates
it in terms of the impossibility that the same
attribute both belongs and does not belong to
the same subject, at the same time and in
the same respect. The qualification by time
and respect is important for distinguishing
genuine contradiction from the attribution
of different predicates under different conditions.
Yet the principle's role in Aristotle's
metaphysics and its defence in Book $\Gamma$
must not be confused with a complete theory
of identity. Logical consistency constrains
how claims about a subject may be made; it
does not by itself determine the full ontological
constitution of that subject. \footnote{See
\textit{Metaphysics} $\Gamma$.3, 1005b19--23,
and $\Gamma$.4, 1006a5--11. Aristotle's defence
is primarily elenctic: he argues that meaningful
assertion already commits the speaker to the
principle, rather than offering an ordinary
demonstration from more basic premises.}

% =============================================================================
% 10. RECIPROCAL ASSESSMENT AND TRANSITION
% =============================================================================

\subsection{Reciprocal Assessment and the Transition to Plotinus}

\phantomsection\label{para:VS.B01.P01.C08.S05.SS10.P001}

Aristotle's principal contribution to the comparative
examination is the development of a differentiated
account of substance and change. The analysis of
categories distinguishes substance from the
predicates and dependent items associated with
it. The later investigation of essence, matter,
form, and composite seeks to explain the constitution
of determinate substances. The account of
potentiality and actuality extends this analysis
to capacities, realization, and the structures
of change. Together, these investigations offer
a substantial framework for examining how
a concrete being is constituted and how it
may undergo different kinds of transformation.

\phantomsection\label{para:VS.B01.P01.C08.S05.SS10.P002}

The strength of the Aristotelian position is
its refusal to treat every change as ontologically
equivalent. It distinguishes alteration from
generation, accidental predicates from essence,
matter from form, and potentiality from actuality.
These distinctions make it possible to ask
what kind of subject persists, what kind of
determination changes, and what explanatory
principle accounts for the unity of a substance.
They also expose the inadequacy of an account
that explains persistence solely through the
continued presence of material constituents.

\phantomsection\label{para:VS.B01.P01.C08.S05.SS10.P003}

The limitations of the comparison are equally
important. Aristotle's priority of substance,
his account of essence, and his hylomorphic
analysis belong to a framework whose central
categories are not identical with the conceptual
foundations of the Philosophy of Viability.
The present ontology cannot adopt substantial
form as a substitute for Identity, treat
the unity of a composite as proof of numerical
singularity, or infer the identity of an
existent solely from the persistence of
its matter. These questions require independent
arguments within the developing system.

\phantomsection\label{para:VS.B01.P01.C08.S05.SS10.P004}

The comparative thesis advanced here is therefore
limited but positive: Aristotle provides a
powerful analysis of substantial determination
and transformation, and his distinctions help
expose the different conditions under which
a subject, a substance, and its constitutive
organization may persist or change. His
framework is a significant interlocutor for
the ontology of identity and manifestation,
but it does not settle the relation between
Being and Identity or the possibility of
plurality of identities. Those questions
remain open within the Philosophy of Viability.

\phantomsection\label{para:VS.B01.P01.C08.S05.SS10.P005}

The next examination turns to Plotinus.
The transition is not simply from a philosophy
of material composition to a philosophy of
immaterial unity. Plotinus develops an account
of reality in which the One, Intellect, and
Soul occupy distinctive explanatory positions,
and in which the relation between unity,
multiplicity, and intelligible determination
becomes central. His position must be examined
on its own terms, especially where it transforms
the questions raised by Plato and Aristotle
concerning substance, form, unity, and the
intelligibility of being. The comparison will
ask what is gained and what is obscured when
the explanation of determinate beings is
related to a principle of unity that is
not itself treated as one being among others.

% =============================================================================
% END OF SECTION 05 OF CHAPTER 08
% =============================================================================